\DocumentMetadata{
  pdfversion=2.0,
  pdfstandard=ua-2,
  lang=en-US,
  tagging=on,
  tagging-setup={math/setup=mathml-SE,tabsorder=structure}
}
\documentclass[11pt,letterpaper]{article}
\usepackage[margin=0.68in,includefoot]{geometry}
\usepackage{newcomputermodern}
\usepackage{amsmath,amscd,mathtools}
\usepackage{tikz,adjustbox}
\providecommand{\square}{\mdlgwhtsquare}
\usepackage[hidelinks]{hyperref}
\hypersetup{
  pdftitle={Algebra First-Year Exam, January 2015, Part 1},
  pdfauthor={University of Florida Department of Mathematics},
  pdfsubject={Graduate mathematics examination},
  pdfdisplaydoctitle=true,
  bookmarksopen=true
}
\setcounter{secnumdepth}{0}
\setlength{\parindent}{0pt}
\setlength{\parskip}{0.28em}
\setlength{\emergencystretch}{3em}
\setlength{\leftmargini}{2.15em}
\setlength{\leftmarginii}{2.65em}
\setlength{\leftmarginiii}{3em}
\pagestyle{empty}
\raggedbottom
\allowdisplaybreaks
\NewTaggingSocketPlug{block/list/label}{examproblem}{%
  \tagstructbegin{tag=Lbl}\tagstructend%
  \tagstructbegin{tag=\LBody}%
  \tagstructbegin{tag=\ProblemHeadingTag,title-o={\ProblemHeadingText}}%
  \tagmcbegin{tag=\ProblemHeadingTag,actualtext={\ProblemHeadingText}}%
  #2%
  \tagmcend%
  \tagstructend%
}
\newenvironment{examproblems}[2]{%
  \def\ProblemHeadingTag{#2}%
  \AssignTaggingSocketPlug{block/list/label}{examproblem}%
  \begin{enumerate}%
  \setcounter{enumi}{#1}%
  \renewcommand{\labelenumi}{\textbf{\theenumi.}}%
  \setlength{\topsep}{0.35em}%
  \setlength{\partopsep}{0pt}%
  \setlength{\itemsep}{0.48em}%
  \setlength{\parsep}{0.18em}%
}{\end{enumerate}}
\newenvironment{examsubparts}{%
  \AssignTaggingSocketPlug{block/list/label}{default}%
  \begin{enumerate}%
  \setlength{\topsep}{0.18em}%
  \setlength{\partopsep}{0pt}%
  \setlength{\itemsep}{0.16em}%
  \setlength{\parsep}{0.1em}%
}{\end{enumerate}}
\newenvironment{examinstructions}{%
  \par\begingroup\itshape
}{%
  \par\endgroup\vspace{0.18em}
}
\makeatletter
\renewcommand\subsection{\@startsection{subsection}{2}{0pt}%
  {-1.25ex plus -.25ex minus -.1ex}%
  {0.55ex plus .1ex}%
  {\normalfont\large\bfseries}}
\renewcommand{\@maketitle}{%
  \newpage\null\vskip -1em%
  {\footnotesize\bfseries
    \href{https://www.ufl.edu/}{University of Florida}, \href{https://math.ufl.edu/}{Department of Mathematics}\par}%
  \vskip -0.5em%
  \tagstructbegin{tag=H1,title={Algebra First-Year Exam, January 2015, Part 1}}%
  \tagpdfparaOff%
  \tagmcbegin{tag=H1}%
    {\LARGE\bfseries\href{https://gma.math.ufl.edu/exams/algebra-first-year/}{Algebra First-Year Exam}, January 2015, Part 1\par}%
  \tagmcend%
  \tagpdfparaOn%
  \tagstructend%
  \vskip 0.25em%
  \tagmcbegin{artifact}%
    \hrule height 1pt%
  \tagmcend%
  \vskip 1em%
}
\makeatother
\title{Algebra First-Year Exam, January 2015, Part 1}
\author{}
\date{}
\begin{document}
\maketitle
\thispagestyle{empty}
\begin{examinstructions}
Answer four problems. If you turn in more than four, only the first four will be graded. Within reason, you may use theorems as long as you state them clearly.
\end{examinstructions}
\begin{examproblems}{0}{H2}
\def\ProblemHeadingText{Problem 1.}
\item
Describe the subgroup
\[
\left\{\begin{pmatrix}a&b\\0&c\end{pmatrix}\;\middle|\;ac\ne 0\right\}\subseteq GL_2(\mathbb{R})
\]
 as a semidirect product \(A\rtimes_\phi B\) with nontrivial \(A,B\). Specify the groups \(A,B\) and the homomorphism \(\phi:B\to\operatorname{Aut}(A)\).
\def\ProblemHeadingText{Problem 2.}
\item
Let~\(G\) be a~\(p\)-group and~\(S\) a finite set on which~\(G\) acts. Denote by \(S^G\subseteq S\) the subset of fixed elements (i.e. \(g(x)=x\) for all \(g\in G\), \(x\in S\)). Show that \(|S|\equiv|S^G|\pmod p\).
\def\ProblemHeadingText{Problem 3.}
\item
Let~\(G\) be a group of order \(3^2\cdot 11\cdot 41\).
\begin{examsubparts}
\item[\textnormal{(i)}]
Show that the 41-Sylow subgroup is normal.
\item[\textnormal{(ii)}]
Show that in fact it is contained in the center.
\item[\textnormal{(iii)}]
Show that~\(G\) has a normal cyclic subgroup of order \(451=11\cdot 41\).
\end{examsubparts}
\def\ProblemHeadingText{Problem 4.}
\item
Suppose~\(p\) is an odd prime.
\begin{examsubparts}
\item[\textnormal{(i)}]
Show that if \(0<a<p\) then the~\(p\)-Sylow subgroup of the symmetric group \(S_{ap}\) is an elementary abelian~\(p\)-group.
\item[\textnormal{(ii)}]
Show that a~\(p\)-Sylow subgroup of \(S_{p^2}\) is nonabelian of order \(p^{p+1}\).
\end{examsubparts}
\def\ProblemHeadingText{Problem 5.}
\item
Suppose~\(G\) is a~\(p\)-group and \(M\le G\) is a maximal subgroup.
\begin{examsubparts}
\item[\textnormal{(i)}]
Show that any maximal subgroup of~\(G\) has index~\(p\).
\item[\textnormal{(ii)}]
Show that~\(M\) contains all commutators and~\(p\)-th powers of elements of~\(G\).
\end{examsubparts}
\def\ProblemHeadingText{Problem 6.}
\item
Let~\(G\) be a finite group and \(H\mathrel{\triangleleft}G\) a normal subgroup. Show that~\(G\) is solvable if and only if~\(H\) and \(G/H\) are solvable.
\end{examproblems}
\end{document}
